%----------------------------------------------------------------------------------------
%       DOCUMENT CONFIGURATION AND INCLUDED PACKAGES
%----------------------------------------------------------------------------------------

\documentclass[english, 10pt,a4paper,oneside,headinclude,footinclude,BCOR5mm]{scrartcl}
\usepackage[nochapters, beramono, eulermath, pdfspacing, dottedtoc ]{classicthesis}

\usepackage{arsclassica}
\usepackage[T1]{fontenc} 
\usepackage[utf8]{inputenc}

\usepackage{graphicx}
\usepackage{tcolorbox}
\usepackage{subfig}
\usepackage{mathtools}
\usepackage{subcaption}
\usepackage{pgfplots}
\usepackage{wrapfig}
\usepackage{tikz}
\usepackage{circuitikz}
\usepackage{siunitx}
\usepackage{tikz-3dplot}
\graphicspath{{figures/}}

\usepackage{enumitem}
\usepackage{lipsum}
\usepackage{varioref}
\usepackage{csquotes}
\usepackage{cancel}

\usepackage{amsmath,amssymb,amsthm}
\usepackage{xfrac}
\usepackage{braket}
\usepackage{upgreek}
\usepackage[version=4]{mhchem}
\usepackage{eulervm}

\usepackage[numbers,sort&compress]{natbib}
\bibliographystyle{unsrtnat}

\hypersetup{
%draft, % Uncomment to remove all links (useful for printing in black and white)
colorlinks=true, breaklinks=true, bookmarks=true,bookmarksnumbered,
urlcolor=webbrown, linkcolor=RoyalBlue, citecolor=webgreen,
pdftitle={},
pdfauthor={\textcopyright},
pdfsubject={},
pdfkeywords={},
pdfcreator={pdfLaTeX},
pdfproducer={LaTeX with hyperref and ClassicThesis}
}


\newtheoremstyle{mytheorem}% Name
  {\topsep}   % Space above
  {\topsep}   % Space below
  {\itshape}  % Body font
  {}          % Indent amount
  {\bfseries} % Head font
  {:}         % Punctuation after theorem head
  {.5em}      % Space after theorem head
  {\thmname{#1}\thmnote{\ (\textbf{#2})}}          % Theorem head spec
\theoremstyle{mytheorem}
\newtheorem{satz}{Satz}[section]


\usepackage[a4paper, margin=2.5cm]{geometry}


\usepackage[english]{babel}

\begin{document}

%----------------------------------------------------------------------------------------
%	CONFIGURE HEADERS, TITLEPAGE, TABLE OF CONTENTS
%----------------------------------------------------------------------------------------
\renewcommand{\sectionmark}[1]{\markright{\spacedlowsmallcaps{#1}}} 
\lehead{\mbox{\llap{\small\thepage\kern1em\color{halfgray} \vline}\color{halfgray}\hspace{0.5em}\rightmark\hfil}} 
\pagestyle{plain} 


\begin{center}
    \large Advanced Computational Physics, University of Strathclyde, WS 2026,\\ Dr. Ben Hourahine

    \vspace{0.5cm}

    {\LARGE\bfseries Assignment 3\par}

    \vspace{0.5cm}

    \large Jörn Wilhelm\\
    Glasgow, 18.\ March 2026

    \vspace{1cm}
\end{center}

\newcommand{\cancelblue}[2]{
  \textcolor{blue}{\cancelto{#1}{\textcolor{black}{#2}}}
}
\newcommand{\overblue}[2]{
  \textcolor{blue}{\overbrace{\textcolor{black}{#2}}^{#1}}
}
\newcommand{\underblue}[2]{
  \textcolor{blue}{\underbrace{\textcolor{black}{#2}}_{#1}}
}
\newcommand{\blue}[1]{\textcolor{blue}{#1}}
\newcommand{\red}[1]{\textcolor{red}{#1}}
%----------------------------------------------------------------------------------------
%	ACTUAL CONTENT
%----------------------------------------------------------------------------------------


\noindent
Running the code yields the three plots of the greens function for points $(0.5,0.5),\,(0.02,0.5)$ and $(0.02,0.02)$\,.
These are shown in figure \ref{fig:greens_functions}.

Furthermore, the results for all possible combinations in task 4, including the $f=0$ case, are printed out, toghether
with the runtime, the standard deviation and the \textit{Gauss-Seidel} result. 
The parameters to set are the grid size $N$ and the number of walkers per core.
Varying the cores from 1 to 16 can be done in the job script.
\\\\
The plot of the speed up factors, i.e. \textit{Amdahl's law},
for varying cores between 1 and 16 is shown in figure \ref{fig:amdahl}. In particular the function 
\begin{align}
  f(n) = \frac{1}{1-p + \frac{p}{n}}
\end{align}
is fitted to the data, yielding 




%----------------------------------------------------------------------------------------
%	BIBLIOGRAPHY
%----------------------------------------------------------------------------------------

%\bibliography{sources}

%----------------------------------------------------------------------------------------
\end{document}